% ===== PRÉAMBULE CANONIQUE — à coller VERBATIM en tête de CHAQUE .tex =====
\documentclass[11pt,a4paper]{article}
\usepackage[utf8]{inputenc}\usepackage[T1]{fontenc}\usepackage{lmodern}
\usepackage[french]{babel}
\usepackage{amsmath,amssymb,mathtools}\usepackage{icomma}
\usepackage{siunitx}\sisetup{locale=FR}  % \qty{}{}, \si{}, \ang{}
\usepackage{xcolor}\usepackage{colortbl}
\usepackage{tikz}\usepackage{pgfplots}\pgfplotsset{compat=1.18}
\usepackage[siunitx,europeanresistors]{circuitikz} % circuits (élec.)
\usepackage[most]{tcolorbox}\usepackage{enumitem}\usepackage{array,booktabs}
\usepackage[a4paper,margin=2.2cm]{geometry}
\usetikzlibrary{arrows.meta,calc,angles,quotes,positioning,patterns,%
  backgrounds,shapes.geometric,decorations.pathmorphing,decorations.markings,intersections}
\definecolor{primary}{RGB}{222,49,99}\definecolor{primarydark}{RGB}{150,22,55}
\definecolor{defcol}{RGB}{0,114,178}\definecolor{methcol}{RGB}{52,92,148}
\definecolor{excol}{RGB}{0,135,90}\definecolor{attncol}{RGB}{200,40,55}
\tcbset{boxbase/.style={enhanced,breakable,boxrule=0.9pt,arc=2.5pt,
  left=3mm,right=3mm,top=2.3mm,bottom=2.3mm,fonttitle=\bfseries,coltitle=white}}
% --- boîtes TUTORIEL (noms EXACTS, ne pas renommer) ---
\newtcolorbox{cle}[1][]{boxbase,colback=defcol!6,colframe=defcol,
  title={\textsc{À retenir}\ifx\relax#1\relax\relax\else\textmd{ : \itshape #1}\fi}}
\newtcolorbox{methode}[1][]{boxbase,colback=methcol!7,colframe=methcol,sharp corners,
  title={\textsc{Méthode}\ifx\relax#1\relax\relax\else\textmd{ --- \itshape #1}\fi}}
\newtcolorbox{exemplebox}[1][]{boxbase,colback=excol!5,colframe=excol,
  title={\textsc{Exemple}\ifx\relax#1\relax\relax\else\textmd{ : \itshape #1}\fi}}
\newtcolorbox{piege}[1][]{boxbase,colback=attncol!6,colframe=attncol,
  title={\textsc{Piège}\ifx\relax#1\relax\relax\else\textmd{ : \itshape #1}\fi}}
\newtcolorbox{remarque}[1][]{boxbase,colback=black!4,colframe=black!45,
  title={\textsc{Remarque}\ifx\relax#1\relax\relax\else\textmd{ : \itshape #1}\fi}}
% --- boîtes EXERCICE / CORRIGÉ (noms EXACTS, auto-comptées) ---
\newtcolorbox[auto counter]{exo}[1][]{boxbase,colback=primary!4,colframe=primary,
  title={\textsc{Exercice}~\thetcbcounter\ifx\relax#1\relax\relax\else\textmd{ --- \itshape #1}\fi}}
\newtcolorbox[auto counter]{corrige}[1][]{boxbase,colback=excol!4,colframe=excol,
  title={\textsc{Corrigé}~\thetcbcounter\ifx\relax#1\relax\relax\else\textmd{ --- \itshape #1}\fi}}
\newcommand{\vv}[1]{\vec{#1}}\newcommand{\ds}{\displaystyle}
\newcommand{\R}{\mathbb{R}}\newcommand{\N}{\mathbb{N}}\newcommand{\Z}{\mathbb{Z}}
% =========================================================================
\begin{document}
\sisetup{per-mode=symbol}

\begin{corrige}[calculer un poids]
On applique $G=mg$.
\begin{enumerate}[label=\alph*)]
  \item $G=10\times\num{9,81}=\qty{98,1}{\newton}$.
  \item $G\approx 10\times 10=\qty{100}{\newton}$.
\end{enumerate}
\end{corrige}

\begin{corrige}[masse ou poids sur la Lune ?]
Le poids $G=mg$ dépend de $g$ ; la masse, elle, est invariable.
\begin{enumerate}[label=\alph*)]
  \item $G_{\text{Terre}}=5\times 10=\qty{50}{\newton}$.
  \item $G_{\text{Lune}}=5\times\num{1,6}=\qty{8}{\newton}$.
  \item Sa masse reste $m=\qty{5}{\kilogram}$ (elle ne change pas de lieu).
\end{enumerate}
\end{corrige}

\begin{corrige}[l'accélération de chute libre]
\begin{enumerate}[label=\alph*)]
  \item Chute libre : $\vv G=m\vv a \Rightarrow a=\dfrac{mg}{m}=g=\qty{9,81}{\metre\per\second\squared}$.
  \item Non : la masse se simplifie, l'accélération vaut $g$ \emph{quelle que soit} la masse.
\end{enumerate}
\end{corrige}

\begin{corrige}[retrouver la masse à partir du poids]
De $G=mg$ : $m=\dfrac{G}{g}=\dfrac{50}{10}=\qty{5}{\kilogram}$.
\end{corrige}

\begin{corrige}[la plume et le marteau]
\begin{enumerate}[label=\alph*)]
  \item Ils touchent le sol \textbf{en même temps}.
  \item En chute libre $a=g$, indépendant de la masse : les deux objets ont la même accélération,
        partent de la même hauteur à vitesse nulle, donc tombent identiquement (expérience d'Apollo 15).
\end{enumerate}
\end{corrige}
\end{document}
